\documentclass[12pt,a4paper]{article}
\usepackage{mathwizards-ingles}
\title{%
  \vspace{-1cm}
  {\Huge\bfseries\textcolor{mes3}{Mes 3 · Semana 11}}\\[0.3cm]
  {\Large Conditionals, Passive Voice e Idioms}\\[0.2cm]
  {\large\color{grissuave} Nivel B1$-$ · Acentos Diversos y Noticias}\\[0.5cm]
  \rule{\textwidth}{1.5pt}
}
\author{}
\date{}

\begin{document}
\maketitle
\thispagestyle{empty}

\begin{cajanivel}
\textbf{Objetivo:} El estudiante se expone a \textbf{acentos globales} (americano, británico, australiano, irlandés) a través de clips de noticias y podcasts variados. Se introducen los \textbf{condicionales} (zero y first), la \textbf{voz pasiva} básica y los primeros \textbf{idioms} (modismos) de alta frecuencia. También se incorporan \textbf{collocations periodísticas} para discutir actualidad.
\end{cajanivel}

% ═══════════════════════════════════════════════════════════════════════════════
\section{Gramática Emergente: Zero Conditional}
% ═══════════════════════════════════════════════════════════════════════════════

\begin{cajagrammar}[Zero Conditional: hechos generales y verdades científicas]
Se usa para expresar situaciones que son \textbf{siempre verdaderas}.

\medskip
\textbf{Estructura:} \eng{If} + Present Simple, Present Simple.

\begin{tabularx}{\textwidth}{@{} X @{}}
\toprule
\textbf{Ejemplos} \\
\midrule
\textit{\textbf{If} you \textbf{heat} water to 100°C, it \textbf{boils}.} \\
\textit{\textbf{If} you \textbf{don't} eat, you \textbf{get} hungry.} \\
\textit{\textbf{If} it \textbf{rains}, the streets \textbf{get} wet.} \\
\textit{\textbf{When} you \textbf{mix} red and blue, you \textbf{get} purple.} \\
\bottomrule
\end{tabularx}

\medskip
\begin{cajaejemplo}[Nota]
En el Zero Conditional, \textbf{if} y \textbf{when} son intercambiables, porque la condición \textit{siempre} se cumple.
\end{cajaejemplo}
\end{cajagrammar}

% ═══════════════════════════════════════════════════════════════════════════════
\section{Gramática Emergente: First Conditional}
% ═══════════════════════════════════════════════════════════════════════════════

\begin{cajagrammar}[First Conditional: situaciones reales y probables en el futuro]
Se usa para hablar de situaciones \textbf{posibles/probables} en el futuro.

\medskip
\textbf{Estructura:} \eng{If} + Present Simple, \eng{will} + verbo base.

\begin{tabularx}{\textwidth}{@{} X @{}}
\toprule
\textbf{Ejemplos} \\
\midrule
\textit{\textbf{If} it \textbf{rains} tomorrow, I\textbf{'ll stay} home.} \\
\textit{\textbf{If} you \textbf{study} hard, you\textbf{'ll pass} the exam.} \\
\textit{\textbf{If} she \textbf{doesn't} hurry, she\textbf{'ll miss} the bus.} \\
\textit{I\textbf{'ll call} you \textbf{if} I \textbf{need} help.} \\
\bottomrule
\end{tabularx}

\medskip
\begin{cajaejemplo}[Error Común]
$\times$ \textit{If it \textbf{will rain}...} \quad $\checkmark$ \textit{If it \textbf{rains}...}\\
Después de \textbf{if}, \textbf{nunca} se usa \textit{will}. La cláusula con ``if'' siempre va en Present Simple.
\end{cajaejemplo}

\medskip
\begin{cajaejemplo}[Contraste Rápido]
\begin{tabularx}{\textwidth}{@{} l X @{}}
\textbf{Zero:} & \textit{If you heat ice, it melts.} (siempre verdadero) \\
\textbf{First:} & \textit{If it rains tomorrow, I'll take an umbrella.} (posible futuro) \\
\end{tabularx}
\end{cajaejemplo}
\end{cajagrammar}

% ═══════════════════════════════════════════════════════════════════════════════
\section{Gramática Emergente: Passive Voice (Introducción)}
% ═══════════════════════════════════════════════════════════════════════════════

\begin{cajagrammar}[Passive Voice: \textit{to be} + past participle]
Se usa cuando el \textbf{objeto de la acción} es más importante que quien la realiza. Muy común en \textbf{noticias y textos formales}.

\medskip
\textbf{Estructura:} Sujeto + \eng{am/is/are/was/were} + \textbf{past participle} (+ by agente).

\medskip
\begin{tabularx}{\textwidth}{@{} l X @{}}
\toprule
\textbf{Activa} & \textbf{Pasiva} \\
\midrule
\textit{They \textbf{built} the house in 1990.} & \textit{The house \textbf{was built} in 1990.} \\
\textit{Millions of people \textbf{speak} English.} & \textit{English \textbf{is spoken} by millions of people.} \\
\textit{They \textbf{will close} the school.} & \textit{The school \textbf{will be closed}.} \\
\textit{Someone \textbf{stole} my phone.} & \textit{My phone \textbf{was stolen}.} \\
\bottomrule
\end{tabularx}

\medskip
\begin{cajaejemplo}[¿Cuándo usar la voz pasiva?]
\begin{itemize}[nosep, leftmargin=0.8cm]
  \item Cuando \textbf{no sabemos} quién hizo la acción: \textit{My wallet \textbf{was stolen}.}
  \item Cuando \textbf{no importa} quién: \textit{The building \textbf{was constructed} in 2005.}
  \item En \textbf{noticias y textos formales}: \textit{Three people \textbf{were injured}.}
\end{itemize}
\end{cajaejemplo}
\end{cajagrammar}

% ═══════════════════════════════════════════════════════════════════════════════
\section{Chunks: Collocations Periodísticas}
% ═══════════════════════════════════════════════════════════════════════════════

\begin{cajachunk}[News Collocations — Para Discutir Noticias y Actualidad]
Estas colocaciones aparecen constantemente en noticias adaptadas (VOA Learning English, BBC Learning English):

\medskip
\begin{tabularx}{\textwidth}{@{} X l @{}}
\eng{break the news} & \esp{dar la noticia} \\
\eng{take action} & \esp{tomar medidas} \\
\eng{raise awareness} & \esp{crear conciencia} \\
\eng{make a statement} & \esp{hacer una declaración} \\
\eng{hold a meeting} & \esp{celebrar una reunión} \\
\eng{reach an agreement} & \esp{llegar a un acuerdo} \\
\eng{face a problem} & \esp{enfrentar un problema} \\
\eng{play a role} & \esp{desempeñar un papel} \\
\eng{have an impact on} & \esp{tener un impacto en} \\
\eng{pay attention to} & \esp{prestar atención a} \\
\eng{according to} & \esp{según} \\
\eng{as a result} & \esp{como resultado} \\
\end{tabularx}
\end{cajachunk}

% ═══════════════════════════════════════════════════════════════════════════════
\section{Chunks: Idioms de Alta Frecuencia}
% ═══════════════════════════════════════════════════════════════════════════════

\begin{cajachunk}[Common Idioms — Modismos que Encontrarás en Podcasts]
Los idioms son expresiones cuyo significado \textbf{no se puede deducir de las palabras individuales}. Se aprenden como \textbf{chunks completos}.

\medskip
\begin{tabularx}{\textwidth}{@{} l X l @{}}
\toprule
\textbf{Idiom} & \textbf{Significado} & \textbf{Equivalente} \\
\midrule
\eng{a piece of cake} & algo muy fácil & \esp{pan comido} \\
\eng{hit the books} & estudiar intensamente & \esp{hincar los codos} \\
\eng{under the weather} & sentirse mal/enfermo & \esp{estar pachuchо} \\
\eng{break the ice} & romper la tensión inicial & \esp{romper el hielo} \\
\eng{it's raining cats and dogs} & llueve muchísimo & \esp{llueve a cántaros} \\
\eng{cost an arm and a leg} & ser muy caro & \esp{costar un ojo de la cara} \\
\eng{let the cat out of the bag} & revelar un secreto & \esp{irse de la lengua} \\
\eng{get cold feet} & acobardarse & \esp{echarse para atrás} \\
\eng{the ball is in your court} & te toca decidir a ti & \esp{la pelota está en tu cancha} \\
\eng{call it a day} & dar por terminado el día & \esp{dar el día por terminado} \\
\eng{once in a blue moon} & muy raramente & \esp{de vez en cuando/de Pascuas a Ramos} \\
\eng{better late than never} & más vale tarde que nunca & \esp{más vale tarde que nunca} \\
\bottomrule
\end{tabularx}
\end{cajachunk}

% ═══════════════════════════════════════════════════════════════════════════════
\section{Chunks: Discutir Acentos y Variantes}
% ═══════════════════════════════════════════════════════════════════════════════

\begin{cajachunk}[Talking About Accents — Frases para las Sesiones de Acentos]
\begin{tabularx}{\textwidth}{@{} X l @{}}
\eng{I find this accent easy/hard to understand.} & \esp{Este acento me parece fácil/difícil.} \\
\eng{She speaks with a British/American accent.} & \esp{Habla con acento británico/americano.} \\
\eng{I'm used to American English.} & \esp{Estoy acostumbrado al inglés americano.} \\
\eng{The pronunciation is very different.} & \esp{La pronunciación es muy diferente.} \\
\eng{Can you speak more slowly, please?} & \esp{¿Puedes hablar más despacio?} \\
\eng{I need to hear that again.} & \esp{Necesito escuchar eso de nuevo.} \\
\end{tabularx}
\end{cajachunk}

% ═══════════════════════════════════════════════════════════════════════════════
\section{Oraciones Listas para Minar}
% ═══════════════════════════════════════════════════════════════════════════════

\begin{cajamineria}[Oraciones \textit{i+1} — Semana 11]

\begin{enumerate}[leftmargin=1.2cm]
  \item \textit{\underline{If} you study every day, you \underline{will} improve quickly.}\\
    {\small\textbf{Objetivo:} estructura del First Conditional (if + present, will + base).}

  \item \textit{If you \underline{heat} ice, it melts. That's a fact.}\\
    {\small\textbf{Objetivo:} Zero Conditional para verdades generales.}

  \item \textit{The building \underline{was constructed} in 1985.}\\
    {\small\textbf{Objetivo:} voz pasiva en pasado (was + past participle).}

  \item \textit{English \underline{is spoken} in many countries around the world.}\\
    {\small\textbf{Objetivo:} voz pasiva en presente (is + past participle).}

  \item \textit{The exam was \underline{a piece of cake}. I finished in 20 minutes.}\\
    {\small\textbf{Objetivo:} el idiom ``a piece of cake'' (muy fácil).}

  \item \textit{I feel \underline{under the weather} today. I think I have a cold.}\\
    {\small\textbf{Objetivo:} el idiom ``under the weather'' (sentirse mal).}

  \item \textit{The government needs to \underline{take action} against pollution.}\\
    {\small\textbf{Objetivo:} la collocation periodística ``take action''.}

  \item \textit{I need to \underline{hit the books} tonight. I have an exam tomorrow.}\\
    {\small\textbf{Objetivo:} el idiom ``hit the books'' (estudiar duro).}

  \item \textit{If she \underline{doesn't} call me, I'll send her a message.}\\
    {\small\textbf{Objetivo:} First Conditional negativo (if + doesn't).}

  \item \textit{\underline{According to} the news, three people were injured.}\\
    {\small\textbf{Objetivo:} la expresión periodística ``according to''.}
\end{enumerate}
\end{cajamineria}

\vfill
\begin{center}
\rule{0.5\textwidth}{0.5pt}\\[0.3cm]
{\small\color{grissuave} Curso de Inmersión en Inglés · Mes 3 · Semana 11 · Nivel B1$-$}
\end{center}

\end{document}
